\section{Introduction}\label{sec:intro}

The Lefschetz operators of a compact hyperk\"ahler manifold generate
much more than the $\mathfrak{sl}_2$ associated with one K\"ahler
class. Together with their lowering operators they generate the
orthogonal Lie algebra of the Beauville--Bogomolov space enlarged by
a hyperbolic plane \cite{verbitsky96,ll97}. For a singular symplectic
variety, intersection cohomology is the natural setting for the same
question. Hard Lefschetz is available there, but an analytic proof of
the additional commutation relations must also specify the domains
of the differential operators near the singularities.

We treat a family of singular moduli spaces on K3 surfaces for which
the regular locus carries a natural hyperk\"ahler metric. The
singularities have positive-dimensional strata and no symplectic
resolution. The main analytic task is to show that the pointwise
hyperk\"ahler identities act on square integrable harmonic forms.
The algebraic task is then to recover all Lefschetz operators from
the three K\"ahler forms of this one metric. We separate these two
arguments: a general algebraic criterion appears in Section
\ref{sec:algebra}, followed by the analysis which verifies it for
the moduli spaces.

\subsection{The moduli spaces}
Let $S$ be a projective K3 surface with $\operatorname{Pic}(S)=\ZZ A$,
where $A$ is ample. Fix a primitive Mukai vector
$v_0=(r_0,\ell A,s_0)$ such that
\begin{equation}\label{eq:family}
 v_0^2=2k,\qquad k\ge2,\qquad r_0\ge2k+1,
 \qquad \gcd(r_0,\ell A^2)=1.
\end{equation}
Write $M=M_A(2v_0)$ and $M_0=M_A(v_0)$ for the moduli spaces of
$A$-Gieseker-semistable sheaves. Under these assumptions the sheaves
are locally free and their slope and Gieseker stability notions
agree as required in Section \ref{sec:geometry}. The singular locus
and its diagonal are
\[
 \Sigma=\operatorname{Sym}^2M_0\subset M,
 \qquad \Delta=\{[F\oplus F]:[F]\in M_0\}\cong M_0.
\]
Their complex dimensions are
\[
 \dim M=8k+2,\qquad \dim\Sigma=4k+4,\qquad \dim\Delta=2k+2.
\]
The space $M$ is an irreducible symplectic variety with $b_2(M)=23$
\cite{pr20,pr23}. Its Beauville--Bogomolov form $q_M$ has real
signature $(3,20)$. It is locally factorial in the Zariski topology
and terminal, and has neither a symplectic resolution nor a smoothing
\cite{kls06,namikawa01,namikawa06}. The qualification on the topology
matters: the transverse analytic singularity at a point of
$\Sigma\setminus\Delta$ is not analytically $\QQ$-factorial.

Put $2n=\dim_{\CC}M$, so $n=4k+1$. Intersection cohomology is graded
in the usual cohomological degrees. A class $x\in H^2(M;\QQ)$ acts
on it by $L_x$. We say that $x$ is a \emph{Lefschetz class} if
\[
 L_x^j:\IH^{2n-j}(M;\QQ)\longrightarrow\IH^{2n+j}(M;\QQ)
\]
is an isomorphism for every $j\ge0$. For such a class there is a
unique operator $\Lambda_x$ of degree $-2$ satisfying
$[L_x,\Lambda_x]=h$, where $h$ acts in degree $j$ by $j-2n$.
The Lie algebra generated by these triples is the
\emph{Looijenga--Lunts--Verbitsky algebra}, abbreviated LLV algebra.

\begin{theorem}\label{thm:main}
For $M$ satisfying \eqref{eq:family}, the Lefschetz classes in
$H^2(M;\QQ)$ are exactly the classes with $q_M(x)\ne0$, and
\[
 \mathfrak g(M;\QQ)\cong
 \mathfrak{so}\big((H^2(M;\QQ),q_M)\oplus U\big),
\]
where $U$ is a hyperbolic plane. Its real form is
$\mathfrak{so}(4,21)$, and its dimension is $300$.
The submodule generated by $1\in\IH^0(M;\CC)$ is the irreducible
representation with highest weight $(n,0,\ldots,0)$ and occurs once.
In particular,
\[
 \operatorname{Sym}^jH^2(M;\CC)\hookrightarrow\IH^{2j}(M;\CC),
 \qquad \dim\IH^{2j}(M)\ge\binom{22+j}{j}\quad(0\le j\le n).
\]
\end{theorem}

Here and below a highest weight is written in orthogonal coordinates
for the root system $B_{12}$. The first example is a quartic K3
surface of Picard rank one with $v_0=(5,A,0)$; then $k=2$ and
$\dim_{\CC}M=18$.

The LLV theorem is known for varieties admitting a symplectic
resolution \cite{fsy22}, and for the isolated-singularity and
orbifold cases treated by Tighe \cite{tighe22}. The spaces above
are outside those classes. Their terminalisation is the space
itself, and the transverse singularity along
$\Sigma\setminus\Delta$ is the closure of the minimal nilpotent
orbit in $\mathfrak{sl}_{2k}$. Its resolution
$T^*\mathbb P^{2k-1}$ is small, so this germ is not analytically
$\QQ$-factorial and cannot be a finite quotient singularity.

\subsection{The analytic statement}
The metric used throughout is the $L^2$ metric on
$M_{\reg}=M\setminus\Sigma$ obtained from Hermitian--Einstein
connections on $S$. Its existence, its K\"ahler class and the
comparison with flat hyperk\"ahler quotients are imported from
\cite{wue15,wue20}; the precise input is stated in Recollection
\ref{thm:imported}. In particular, the metric comparison imported
there is a bound on the metric itself, without a derivative bound.
The first-order estimates needed below are proved separately.

Let $I_a$, $\omega_a$, $a=1,2,3$, denote the parallel complex
structures and K\"ahler forms of this metric. On differential forms
let $W_a$ be the Weil operator, acting on $(p,q)$-forms for $I_a$ by
$i^{p-q}$, and put $d_a^c=W_a^{-1}dW_a$. All closed extensions below
are taken in the $L^2$ space of the regular locus.

\begin{theorem}\label{thm:analytic}
For the metric $g$ on $M_{\reg}$:
\begin{enumerate}
\item $d_{\min}=d_{\max}$, and the maximal de Rham complex has
finite-dimensional cohomology and closed range. There is a canonical
isomorphism
\[
 \IH^\bullet(M;\RR)\cong H^\bullet_{(2)}(M_{\reg},g)
 \cong\cH^\bullet,
 \qquad \cH=\ker d_{\max}\cap\ker\delta_{\max}.
\]
\item The isomorphism intertwines the $H^\bullet(M;\RR)$ action
with wedge product by bounded closed representatives. It identifies
the intersection pairing with integration of wedge products.
\item Each $W_a$ preserves the Gaffney domain
$V=\Dom d_{\max}\cap\Dom\delta_{\max}$. Each harmonic form is
$d_a^c$-closed and $d_a^c$-coclosed. Wedge product by $\omega_a$ and
its pointwise adjoint preserve $\cH$ and satisfy Verbitsky's
$\mathfrak{so}(4,1)$ relations there.
\end{enumerate}
\end{theorem}

The first assertion follows from local cone cohomology and the
characterisation of the intersection complex. The third is more
sensitive to the metric. Equality of the formal Laplacians for $d$
and $d_a^c$ on the regular locus does not identify their closed
extensions. It is the invariance of $V$ that permits the K\"ahler
identities to be used on harmonic forms.

Two features of the proof are useful beyond this family. First, the
K\"ahler cone theorem is extended from smooth links to singular
links with discrete Hodge spectrum and the $L^2$ Stokes property.
The potentially problematic modes are homogeneous closed and
coclosed forms at the critical indicial exponents. The K\"ahler
identities force their $d^c$ derivatives into gaps in the exponents
of adjacent degrees; an integration-by-parts argument on the link
treats the endpoint. This is the mechanism of Cheeger's smooth-link
theorem \cite{cheeger83}, arranged here to permit induction on depth.

Second, the comparison of Gaffney domains keeps track of directions.
At a single stratum, a metric error of order the distance to that
stratum, with bounded first derivative, is sufficient. Near the
diagonal there are two distances. Errors in the normal and mixed
directions must vanish at the nearer stratum, whereas tangential
errors may be of the larger order. Tangential regularity on a
product model controls the latter. These estimates are verified
using finite-dimensional slices of the gauge-theoretic quotient.

\subsection{From one metric to all classes}
Theorem \ref{thm:analytic} provides Lefschetz triples only on the
three-plane
$\mathcal P=\langle[\omega_1],[\omega_2],[\omega_3]\rangle$.
The Beauville--Bogomolov form is a positive multiple of the Euclidean
form on this plane. The monodromy of locally trivial deformations
has Zariski-dense image in the special orthogonal group of $H^2$.
The operator relations needed for the LLV theorem are algebraic in
at most three classes. Every nondegenerate three-plane is an
orthogonal translate of $\mathcal P_{\CC}$; monodromy invariance
therefore extends those relations to all classes. This proves
Theorem \ref{thm:main} without requiring an analytic theorem for
every K\"ahler class.

Section \ref{sec:monodromy} also shows that a subgroup of finite
index of monodromy acts through the spin group of $H^2$. The
intersection pairing and the positive $L^2$ inner product make the
possible centralising factor finite. Section
\ref{sec:representations} applies this to the eighteen-dimensional
example. Its complete decomposition requires a specified
isomorphism of graded rational Hodge structures motivated by the
BPS and $\chi$-independence results of
\cite{dhsm22,dhksv26,gwz26}. We state that premise explicitly; a
calculation of graded dimensions alone does not identify an LLV
representation.

The analytic background is the $L^2$ theory of Cheeger and
Cheeger--Goresky--MacPherson \cite{cheeger80,cgm82}, its Hilbert
complex formulation \cite{brl92,brl93}, and the edge theories of
\cite{hm05,almp12,almp18}. Our metric is not assumed to have an
iterated edge normal form. The quotient stratification is consistent
with \cite{ds97,mayrand22}; the charts used here additionally
control the metric and all three complex structures. No
identification of the harmonic type decomposition with Saito's
Hodge structure is used in the LLV theorem or in the conditional
representation argument.
